%
% effort-comparison.tex
%
% Development Effort and Investment Comparison:
% xboing-python vs xboing-c
%
% Compile: pdflatex effort-comparison.tex
%

\documentclass[a4paper,10pt]{article}
\usepackage[margin=1in]{geometry}
\usepackage{booktabs}
\usepackage{listings}
\usepackage{xcolor}
\usepackage{enumitem}
\usepackage{longtable}
\usepackage[colorlinks=true,linkcolor=blue,citecolor=blue,urlcolor=blue]{hyperref}

\lstset{
  basicstyle=\ttfamily\small,
  keywordstyle=\color{blue!80!black},
  commentstyle=\color{green!50!black},
  stringstyle=\color{red!60!black},
  breaklines=true,
  frame=single,
  framerule=0.3pt,
  rulecolor=\color{gray!50},
  backgroundcolor=\color{gray!5},
  columns=fullflexible
}

\newcommand{\repo}[1]{\texttt{#1}}
\newcommand{\model}[1]{\textbf{#1}}

\begin{document}

\title{Development Effort and Investment Comparison:\\
  \texttt{xboing-python} vs \texttt{xboing-c}}
\author{XBoing Modernization Project}
\date{June 2026 --- Generated from git history}
\hypersetup{
  pdftitle={XBoing Development Effort Comparison},
  pdfauthor={Jim Freeman},
  pdfsubject={Effort Analysis},
  pdfcreator={pdflatex}
}
\maketitle

\tableofcontents
\newpage

% ===================================================================
\section{Overview}
% ===================================================================

This report compares the development effort invested in two ports of
the classic XBoing (1993, Justin C.\ Kibell) breakout game:

\begin{itemize}
  \item \repo{xboing-python}
    (\url{https://github.com/jmf-pobox/xboing-python}) ---
    Python~3.10+ / pygame.
  \item \repo{xboing-c}
    (\url{https://github.com/jmf-pobox/xboing-c}) ---
    C11 / SDL2.
\end{itemize}

Both projects share the same maintainer and the same goal: a faithful
port of the original XBoing.  They differ in language, timing,
tooling, and the degree of AI-assisted development.  This report
quantifies those differences using git history as the primary data
source and correlates development phases with the Claude model
generations available at each point.

% ===================================================================
\section{Timeline}
% ===================================================================

\begin{center}
\begin{tabular}{@{}lll@{}}
\toprule
& \textbf{xboing-python} & \textbf{xboing-c} \\
\midrule
First commit & 2025-05-03 & 2025-11-26 \\
Latest commit & 2026-06-02 & 2026-05-30 \\
Calendar span & 13 months & 6 months \\
Active development days & 24 & 30 \\
Elapsed active weeks & 6 & 16 \\
\bottomrule
\end{tabular}
\end{center}

\paragraph{Activity pattern.}
\repo{xboing-python} had a single intense burst in May~2025
(301~commits in 4~weeks), then went dormant for 6~months, with brief
activity in November~2025 (5~commits) and June~2026 (3~commits).
\repo{xboing-c} started 7~months later but sustained development
across 16~weeks from February--May~2026, with a peak of 162~commits
in February alone.

% ===================================================================
\section{Volume and Velocity}
% ===================================================================

\subsection{Aggregate Metrics}

\begin{center}
\begin{tabular}{@{}lrr@{}}
\toprule
\textbf{Metric} & \textbf{xboing-python} & \textbf{xboing-c} \\
\midrule
Total commits & 310 & 264 \\
Merged pull requests & 18 & 137 \\
Lines added & 64,568 & 142,861 \\
Lines removed & 21,066 & 22,933 \\
Net lines & +43,502 & +119,928 \\
Source lines (current) & 12,372 & 23,332 \\
Test lines (current) & 5,674 & 25,442 \\
Test functions & 202 & 1,115 \\
Source files & 82 (.py) & 43 (.c) \\
Test files & 40 & 59 \\
\bottomrule
\end{tabular}
\end{center}

\paragraph{Observation.}  \repo{xboing-python} has more raw commits
(310 vs 264) but far fewer PRs (18 vs 137).  The Python repo's early
history consists of direct pushes to master; the C~repo adopted a
strict PR workflow from the start, with every change going through
review.  The C~repo's net line output is 2.75$\times$ higher, and its
test corpus is 4.5$\times$ larger.

\subsection{Monthly Breakdown}

\begin{longtable}{@{}llrrr@{}}
\toprule
\textbf{Month} & \textbf{Repo} & \textbf{Commits} &
  \textbf{Lines +} & \textbf{Lines --} \\
\midrule
\endfirsthead
\toprule
\textbf{Month} & \textbf{Repo} & \textbf{Commits} &
  \textbf{Lines +} & \textbf{Lines --} \\
\midrule
\endhead
2025-05 & xboing-python & 301 & 42,817 & 17,129 \\
2025-06 & xboing-python & 1 & 7 & 651 \\
\midrule
2025-11 & xboing-python & 5 & 17,421 & 2,024 \\
2025-11 & xboing-c & 2 & 39,541 & 39 \\
\midrule
2026-02 & xboing-c & 162 & 64,778 & 18,366 \\
2026-03 & xboing-c & 52 & 8,855 & 773 \\
2026-04 & xboing-c & 8 & 5,229 & 889 \\
2026-05 & xboing-c & 40 & 24,458 & 2,866 \\
\midrule
2026-06 & xboing-python & 3 & 4,323 & 1,262 \\
\bottomrule
\end{longtable}

The two projects never had concurrent active development.
\repo{xboing-python}'s burst ended in May~2025; \repo{xboing-c}'s
began in November~2025.  The February~2026 spike in \repo{xboing-c}
(162~commits, +64K lines) represents the SDL2 rewrite of the core
game systems.

% ===================================================================
\section{AI-Assisted Development}
% ===================================================================

Both projects are entirely AI-written.  \repo{xboing-c} was written with Claude Code (Opus~4/4.5/4.6).
\repo{xboing-python} was written with Cursor.  In May~2025,
Cursor's available models were Claude~3.5~Sonnet~v2 (the default
coding model) and GPT-4o.  The difference between the two projects
is in model generation, tooling maturity, and workflow conventions.

\subsection{Co-Authorship Tags}

Commits containing a \texttt{Co-Authored-By: Claude} tag or
\texttt{Generated with Claude Code} marker:

\begin{center}
\begin{tabular}{@{}lrrr@{}}
\toprule
\textbf{Repo} & \textbf{Total} & \textbf{Tagged} &
  \textbf{Tag Rate} \\
\midrule
xboing-python & 310 & 5 & 1.6\% \\
xboing-c & 264 & 202 & 76.5\% \\
\bottomrule
\end{tabular}
\end{center}

\paragraph{Interpretation.}  The low tag rate on
\repo{xboing-python} reflects the use of Cursor, which does not add
co-authorship tags.  The 5~tagged commits (November~2025 and
June~2026) correspond to the few occasions where Claude Code was
used on the Python repo.  \repo{xboing-c} used Claude Code
exclusively, where tagging is standard.

\subsection{Tagged Commits by Month}

\begin{center}
\begin{tabular}{@{}lrr@{}}
\toprule
\textbf{Month} & \textbf{xboing-c (tagged)} &
  \textbf{xboing-python (tagged)} \\
\midrule
2025-05 & --- & 0 of 301 \\
2025-11 & 1 & 2 \\
2026-02 & 82 & --- \\
2026-03 & 44 & --- \\
2026-04 & 33 & --- \\
2026-05 & 42 & --- \\
2026-06 & --- & 3 \\
\bottomrule
\end{tabular}
\end{center}

% ===================================================================
\section{Claude Model Availability Windows}
% ===================================================================

The following table maps Claude model releases against the development
timelines of both projects.  Release dates are approximate based on
public announcements.

\begin{longtable}{@{}llp{6cm}@{}}
\toprule
\textbf{Date} & \textbf{Model} & \textbf{Project Activity} \\
\midrule
\endfirsthead
\toprule
\textbf{Date} & \textbf{Model} & \textbf{Project Activity} \\
\midrule
\endhead
2024-06 & \model{Sonnet 3.5} &
  Neither project started. \\
2024-10 & \model{Sonnet 3.5 v2}, \model{Haiku 3.5} &
  Neither project started. \\
2025-02 & \model{Sonnet 3.5 v2} (Claude Code early access) &
  Neither project started. Claude Code in limited
  preview. \\
\midrule
2025-04 & \model{Sonnet 3.5 v2} (Claude Code GA) &
  \repo{xboing-python} starts (May~3) using Cursor.
  301~commits in May.  Cursor's models: Claude~3.5
  Sonnet~v2 (default) and GPT-4o. \\
2025-06 & \model{Sonnet 4 (preview)} &
  \repo{xboing-python} goes dormant (1~commit in June). \\
2025-10 & \model{Sonnet 3.6} &
  Both projects dormant. \\
\midrule
2025-11 & \model{Sonnet 3.6} &
  \repo{xboing-c} baseline commit (Nov~26).
  \repo{xboing-python} brief activity (5~commits). \\
2025-12 & \model{Opus 4} &
  Both projects dormant.  Opus~4 is the first model
  with strong agentic coding capability. \\
2026-01 & \model{Sonnet 4.5}, \model{Haiku 4.5} &
  \repo{xboing-c} dormant between baseline and main
  development phase. \\
\midrule
2026-02 & \model{Opus 4} (dominant coding model) &
  \repo{xboing-c} primary development burst begins.
  162~commits, 82~Claude co-authored.  The SDL2
  rewrite of core game systems happens under Opus~4. \\
2026-03 & \model{Opus 4.5} &
  \repo{xboing-c} continues.  52~commits,
  44~Claude co-authored. \\
2026-04 & \model{Opus 4.5} &
  \repo{xboing-c} lower intensity.  8~commits,
  33~Claude co-authored (includes branches). \\
2026-05 & \model{Opus 4.6}, \model{Sonnet 4.6} &
  \repo{xboing-c} second surge.  40~commits,
  42~Claude co-authored.  Savegame system, dialogue,
  visual polish. \\
2026-06 & \model{Opus 4.6} &
  \repo{xboing-python} brief activity (3~commits,
  dependency upgrades + PR cleanup). \\
\bottomrule
\end{longtable}

\subsection{Key Timing Observations}

\begin{enumerate}
  \item \textbf{Both projects are entirely AI-written.}  The
    difference is not whether AI was used, but which model generation
    and tooling were available.

  \item \textbf{\repo{xboing-python} was built with Cursor and
    Sonnet~3.5~v2 / GPT-4o.}  The May~2025 burst used Cursor with
    Claude~3.5~Sonnet~v2 as the primary coding model.  Cursor in
    this period offered inline completion, chat, and composer ---
    but not the agentic autonomy of later Claude Code (no
    subagents, no PR automation, no hook-based workflows).

  \item \textbf{\repo{xboing-c} was built with Opus~4/4.5/4.6
    and mature Claude Code.}  The entire SDL2 rewrite
    (Feb--May~2026) coincides with the most capable Claude models
    to date and a substantially more mature tooling ecosystem.

  \item \textbf{Model capability inflection.}  The jump from
    Sonnet~3.5~v2 (May~2025) to Opus~4+ (Feb--May~2026) represents
    a step change in coding capability: longer context, better
    tool use, stronger multi-file reasoning, and more reliable
    code generation.  This capability difference directly impacts
    output volume and quality per session.

  \item \textbf{Tooling paradigm shift.}  Beyond model capability,
    the tooling paradigm changed between the two projects.  Cursor
    (May~2025) is an IDE with AI-assisted editing --- the developer
    drives, the AI assists.  Claude Code (Feb--May~2026) is an
    agentic CLI where the AI drives autonomously --- creating
    branches, running tests, reviewing PRs, resolving comments,
    and merging.  Features like automated PR review loops, merge
    automation, subagents, and multi-agent orchestration were
    available only for \repo{xboing-c}.
\end{enumerate}

% ===================================================================
\section{Investment Comparison}
% ===================================================================

\subsection{Effort Density}

\begin{center}
\begin{tabular}{@{}lrr@{}}
\toprule
\textbf{Metric} & \textbf{xboing-python} & \textbf{xboing-c} \\
\midrule
Active development days & 24 & 30 \\
Commits per active day & 12.9 & 8.8 \\
Lines added per active day & 2,690 & 4,762 \\
PRs per active day & 0.75 & 4.6 \\
Test functions per active day & 8.4 & 37.2 \\
\bottomrule
\end{tabular}
\end{center}

\paragraph{Interpretation.}  \repo{xboing-c} produces 1.8$\times$
more lines per active day and 4.4$\times$ more test functions per
active day.  Both projects are entirely AI-written, so the
productivity difference reflects the model and tooling generations:
Opus~4+ with agentic Claude Code (subagents, PR automation, hooks)
vs Sonnet~3.5~v2 / GPT-4o with Cursor (assisted editing).

\subsection{Quality Investment}

\begin{center}
\begin{tabular}{@{}lrr@{}}
\toprule
\textbf{Metric} & \textbf{xboing-python} & \textbf{xboing-c} \\
\midrule
Test:source line ratio & 0.46:1 & 1.09:1 \\
Test functions & 202 & 1,115 \\
Fuzz targets & 0 & 4 \\
Golden screenshot tests & No & Yes \\
Integration test suites & 1 & 6 \\
Merged PRs (review gates) & 18 & 137 \\
\bottomrule
\end{tabular}
\end{center}

\repo{xboing-c} has more than double the test-to-source ratio and
5.5$\times$ more test functions.  It includes fuzz testing for all
I/O parsers and golden screenshot tests for visual fidelity.  Every
change goes through a PR with automated CI and reviewer sign-off
(137~merged PRs vs 18).

\subsection{Feature Output}

\begin{center}
\begin{tabular}{@{}lrr@{}}
\toprule
\textbf{Metric} & \textbf{xboing-python} & \textbf{xboing-c} \\
\midrule
Game modes implemented & 4 of 16 & 16 of 16 \\
Block types with behavior & $\sim$12 of 30 & 30 of 30 \\
Visual effects & 0 of 5 & 5 of 5 \\
Persistence systems & 0 of 3 & 3 of 3 \\
Specials with behavior & 2 of 8 & 8 of 8 \\
\bottomrule
\end{tabular}
\end{center}

% ===================================================================
\section{Synthesis}
% ===================================================================

\begin{enumerate}
  \item \textbf{Both projects are entirely AI-written; the model
    and tooling generation is the variable.}  The output difference
    --- 2.75$\times$ more code, 5.5$\times$ more tests, 4$\times$
    more game modes --- is not explained by one project using AI
    and the other not.  Both are AI-written.  The difference is
    Cursor + Sonnet~3.5~v2 (May~2025) vs Claude Code +
    Opus~4/4.5/4.6 (Feb--May~2026).

  \item \textbf{Active development days are comparable.}
    \repo{xboing-python}: 24~active days.  \repo{xboing-c}:
    30~active days.  The 25\% difference in active days does not
    explain the 2--5$\times$ difference in output.  The multiplier
    is model capability and tooling maturity.

  \item \textbf{Agentic vs assisted is the tooling divide.}
    Cursor (2025) assists a developer who drives the workflow.
    Claude Code (2026) drives autonomously --- creating branches,
    writing code, running quality gates, opening PRs, addressing
    review comments, and merging.  The 137~merged PRs on
    \repo{xboing-c} vs 18 on \repo{xboing-python} reflect this
    paradigm difference: automated PR creation, review loops, and
    merge automation were not part of the Cursor workflow.

  \item \textbf{The PR workflow gap is a consequence of tooling,
    not discipline.}  \repo{xboing-python} pushed directly to
    master (301~commits in May~2025) because Cursor does not
    automate branch/PR/merge workflows.  \repo{xboing-c} used PRs
    for every change because Claude Code automates the entire
    cycle.  The later project's higher test coverage and review
    discipline follow from the tooling.

  \item \textbf{The Python port's foundations are sound.}  Despite
    the completeness gap, \repo{xboing-python} has a clean MVC
    architecture, strict type checking (mypy + pyright in strict
    mode), protocol-based design, and dependency injection.
    Bringing the Python port to parity would require implementing
    the missing game systems, not rearchitecting what exists.
\end{enumerate}

\end{document}
